%---------------------------------------------------------------------
% Part three: the Prophet's multifaceted role. Four questions.
%
% This "بحیثیت ..." format is the signature question type of this
% course -- it appeared, in one label or another, in all four papers.
%---------------------------------------------------------------------

\section*{حصہ سوم --- نبی کریم \saw کی کثیر الجہتی شخصیت}
\markboth{کثیر الجہتی شخصیت}{}

\begin{nukta}[یہ سوال ہر بار آتا ہے --- سانچہ یاد رکھیے]
پرچے میں ''بحیثیت ...'' کے عنوان سے ہر بار ایک نیا پہلو پوچھا جاتا ہے ---
معلم، مدبر، قانون ساز، سربراہِ خاندان، محسنِ انسانیت، مبلغ۔ ہر جواب کا
سانچہ ایک جیسا رکھیے: \textbf{(۱)} تعارف، \textbf{(۲)} دو تین ٹھوس مثالیں
واقعات یا احکام کی صورت میں، \textbf{(۳)} متعلقہ حدیث یا آیت، \textbf{(۴)}
خلاصہ۔
\end{nukta}

%---------------------------------------------------------------------
\begin{sawal}{سوال \textenglish{13} \ \textnormal{\small(بہار \textenglish{2013}، سوال \textenglish{6} اور خزاں \textenglish{2017}، سوال \textenglish{7})}}
نبی کریم \saw کی ''معلمِ انسانیت'' اور ''معلمِ اخلاق'' بحیثیت پر مفصل نوٹ
قلمبند کریں۔
\end{sawal}

\begin{hal}
\textbf{تعارف۔} نبی کریم \saw نے خود اپنی بعثت کا مقصد بیان فرمایا:
\ar{إِنَّمَا بُعِثْتُ لِأُتَمِّمَ صَالِحَ الْأَخْلَاقِ} --- مجھے اسی لیے
بھیجا گیا کہ نیک اخلاق کی تکمیل کروں (موطأ امام مالک)۔

\medskip
\textbf{بحیثیتِ معلم --- طریقۂ تعلیم:}
\begin{itemize}\setlength{\itemsep}{1pt}
  \item \textbf{عملی نمونہ} --- زبانی نصیحت سے زیادہ اپنے عمل سے سکھانا؛
        اسی لیے فرمایا: ''تم میں سب سے بہتر وہ ہے جو قرآن سیکھے اور
        سکھائے'' (بخاری)۔
  \item \textbf{تدریج} --- احکام (نماز، شراب کی حرمت) مرحلہ وار نازل ہوئے،
        یکبارگی نہیں۔
  \item \textbf{سوال کا خیر مقدم} --- صحابہؓ کے سوالات کو ناپسند نہیں
        کیا جاتا تھا؛ حدیثِ جبریل اسی کی مثال ہے، جس میں جبریل\,علیہ السلام
        نے سوال پوچھ کر صحابہؓ کو ایمان، اسلام اور احسان کی تعلیم دلوائی۔
  \item \textbf{مساوات} --- امیر و غریب، آزاد و غلام سب کو ایک ہی مجلس میں
        بٹھایا۔
\end{itemize}

\medskip
\textbf{بحیثیتِ معلمِ اخلاق --- بنیادی تعلیمات:} سچائی، امانت، وعدہ نبھانا،
عفو و درگزر، حسنِ سلوک، اور تکبر سے اجتناب۔ حضرت انسؓ فرماتے ہیں کہ نبی
کریم \saw دس سال ان کی خدمت میں رہے، کبھی ''اُف'' تک نہ کہا (بخاری) ---
اخلاق کی یہ بلندی محض نصیحت سے نہیں، عملی نمونے سے سکھائی گئی۔

\medskip
\textbf{خلاصہ۔} نبی کریم \saw کی تعلیم کا مرکز \emph{کردار سازی} تھا ---
علم کو عمل اور اخلاق سے جوڑنا، جو آج بھی تعلیم کے میدان میں رہنما اصول
ہے۔
\end{hal}

%---------------------------------------------------------------------
\begin{sawal}{سوال \textenglish{14} \ \textnormal{\small(بہار \textenglish{2013}، سوال \textenglish{7--8})}}
نبی کریم \saw کی ''مدبر و منظم'' اور ''قانون ساز'' بحیثیت پر مفصل نوٹ
قلمبند کریں۔
\end{sawal}

\begin{hal}
\textbf{بحیثیتِ مدبر و منظم۔} مدینہ پہنچتے ہی نبی کریم \saw نے ایک بکھری
ہوئی آبادی کو منظم ریاست میں بدل دیا:
\begin{itemize}\setlength{\itemsep}{1pt}
  \item \textbf{مسجدِ نبوی کی تعمیر} --- عبادت کے ساتھ ساتھ حکومتی امور،
        تعلیم اور مہمانوں کے قیام کا مرکز۔
  \item \textbf{عقدِ مواخات} --- مہاجرین و انصار کے درمیان بھائی چارہ
        (تفصیل سوال \textenglish{7})۔
  \item \textbf{میثاقِ مدینہ} --- تحریری آئین کے ذریعے مختلف اقوام کو ایک
        ریاست میں پرو دیا (تفصیل سوال \textenglish{7})۔
  \item \textbf{انتظامی تقسیم} --- جنگوں میں سالار مقرر کرنا، صدقات و
        زکوٰۃ جمع کرنے کے لیے عاملین بھیجنا، اور مشاورت (شوریٰ) کا اصول
        اپنانا: \ar{وَشَاوِرْهُمْ فِي الْأَمْرِ} (آل عمران،
        \textenglish{159})۔
\end{itemize}

\medskip
\textbf{بحیثیتِ قانون ساز۔} نبی کریم \saw نے قرآن کی روشنی میں ایک مکمل
قانونی نظام قائم کیا:
\begin{itemize}\setlength{\itemsep}{1pt}
  \item \textbf{فوجداری قانون} --- قصاص و دیت کے احکام، حدود کا نفاذ ---
        سب کے لیے یکساں، حتیٰ کہ فرمایا: ''اگر فاطمہ بنتِ محمد بھی چوری
        کرتی تو میں اس کا ہاتھ کاٹ دیتا'' (بخاری و مسلم)۔
  \item \textbf{خاندانی قانون} --- نکاح، طلاق، وراثت کے مفصل احکام۔
  \item \textbf{معاہدات و بین الاقوامی قانون} --- صلح حدیبیہ اور یہود سے
        معاہدے اسی کی مثالیں۔
  \item \textbf{عدالتی نظام} --- تنازعات کا فیصلہ خود کرتے، گواہوں اور
        قسم کا اصول رائج کیا۔
\end{itemize}

\medskip
\textbf{خلاصہ۔} دس سالہ مدنی دور میں نبی کریم \saw نے محض ایک مذہبی رہنما
نہیں بلکہ ایک مکمل ریاست کے سربراہ، منتظم اور قانون ساز کی حیثیت سے کام
کیا --- جدید اصطلاح میں یہ ''ریاستِ مدینہ'' کہلاتی ہے۔
\end{hal}

%---------------------------------------------------------------------
\begin{sawal}{سوال \textenglish{15} \ \textnormal{\small(خزاں \textenglish{2012}، سوال \textenglish{5} اور بہار \textenglish{2022}، سوال \textenglish{7})}}
نبی کریم \saw کی ''سربراہِ خاندان'' بحیثیت اور ''رحمۃ للعالمین'' پر مفصل
نوٹ لکھیں۔
\end{sawal}

\begin{hal}
\textbf{بحیثیتِ سربراہِ خاندان۔} نبی کریم \saw گھریلو زندگی میں بھی مکمل
نمونہ تھے:
\begin{itemize}\setlength{\itemsep}{1pt}
  \item \textbf{ازواجِ مطہراتؓ سے حسنِ سلوک} --- ''تم میں سب سے بہتر وہ
        ہے جو اپنے اہل کے لیے بہتر ہو، اور میں اپنے اہل کے لیے تم میں سب
        سے بہتر ہوں'' (ترمذی)۔
  \item \textbf{گھر کے کام میں شرکت} --- اپنے کپڑے خود سیتے، جوتا خود
        گانٹھتے، گھر کا کام کرتے (بخاری)۔
  \item \textbf{اولاد سے محبت} --- حضرت فاطمہؓ کے آنے پر خود کھڑے ہو کر
        استقبال کرتے، اور نواسوں حسنؓ و حسینؓ کو کندھوں پر بٹھاتے۔
  \item \textbf{مشاورت اور نرمی} --- گھریلو فیصلوں میں سختی کے بجائے مشورہ
        اور نرمی کا رویہ۔
\end{itemize}

\medskip
\textbf{رحمۃ للعالمین۔} قرآن نے نبی کریم \saw کی بعثت کا مقصد ہی رحمت
قرار دیا: \ar{وَمَا أَرْسَلْنَاكَ إِلَّا رَحْمَةً لِّلْعَالَمِينَ}
(الانبیاء، \textenglish{107})۔ یہ رحمت مسلمانوں تک محدود نہ تھی:

\begin{itemize}\setlength{\itemsep}{1pt}
  \item \textbf{فتح مکہ کی عام معافی} --- سخت ترین دشمنوں کو بھی معاف کر
        دیا۔
  \item \textbf{غیر مسلموں سے حسنِ سلوک} --- طائف میں شدید زخمی ہونے کے
        باوجود بددعا نہ دی، بلکہ ہدایت کی دعا کی۔
  \item \textbf{جانوروں اور فطرت پر رحم} --- بلا وجہ درخت کاٹنے یا جانور
        کو ستانے سے منع فرمایا۔
  \item \textbf{عالمگیریت} --- پیغام کسی ایک قوم کے لیے نہیں: \ar{وَمَا
        أَرْسَلْنَاكَ إِلَّا كَافَّةً لِّلنَّاسِ} (سبا، \textenglish{28})۔
\end{itemize}

\medskip
\textbf{خلاصہ۔} گھر میں نرمی اور باہر رحمت --- دونوں ایک ہی سیرت کے دو
پہلو ہیں۔
\end{hal}

%---------------------------------------------------------------------
\begin{sawal}{سوال \textenglish{16} \ \textnormal{\small(خزاں \textenglish{2012}، سوال \textenglish{8}؛ بہار \textenglish{2022}، سوال \textenglish{6}؛ خزاں \textenglish{2017}، سوال \textenglish{8})}}
نبی کریم \saw کی معاشی زندگی اور معاشی تعلیمات پر جامع نوٹ تحریر کریں۔
\end{sawal}

\begin{hal}
یہ موضوع \textbf{تین پرچوں میں} آیا ہے --- سب سے زیادہ دہرایا جانے والا
موضوعاتی سوال۔

\medskip
\textbf{ذاتی معاشی زندگی۔} نبی کریم \saw نے جوانی میں تجارت کی، خصوصاً
حضرت خدیجہؓ کا مالِ تجارت لے کر شام گئے --- دیانت اور امانت کی وجہ سے
\textbf{''صادق و امین''} کا لقب پایا۔ نبوت کے بعد سادہ زندگی اپنائی؛ گھر
میں کئی کئی دن چولہا نہ جلتا، مگر کبھی کسی سے سوال نہ کیا۔ حدیث میں ہے:
''رسول اللہ \saw دنیا کی سب سے زیادہ سخی ہستی تھے، اور رمضان میں آپ کی
سخاوت تیز ہوا سے بھی زیادہ ہوتی'' (بخاری) --- \ar{أَجْوَدُ بِالْخَيْرِ
مِنَ الرِّيحِ الْمُرْسَلَةِ}۔

\medskip
\textbf{معاشی تعلیمات:}
\begin{itemize}\setlength{\itemsep}{1pt}
  \item \textbf{سود کی حرمت} --- \ar{وَأَحَلَّ اللَّهُ الْبَيْعَ وَحَرَّمَ
        الرِّبَا} (البقرہ، \textenglish{275})۔
  \item \textbf{زکوٰۃ کا نظام} --- مال کو گردش میں رکھنے اور دولت کو چند
        ہاتھوں تک محدود نہ ہونے دینے کا اصول۔
  \item \textbf{دیانتِ تجارت} --- ناپ تول میں کمی، ملاوٹ اور دھوکے کی
        سخت ممانعت: جو ہمیں دھوکا دے، وہ ہم میں سے نہیں (مسلم)۔
  \item \textbf{مزدور کے حقوق} --- ''مزدور کی مزدوری اس کا پسینہ خشک ہونے
        سے پہلے ادا کرو'' (ابن ماجہ)۔
  \item \textbf{بیت المال کا قیام} --- عوامی وسائل کی شفاف تقسیم کی
        بنیاد۔
\end{itemize}

\medskip
\textbf{خلاصہ۔} نبی کریم \saw کی معاشی تعلیمات کا مرکزی اصول \textbf{عدل
اور دیانت} ہے --- نہ سرمایہ کی اجارہ داری، نہ استحصال، بلکہ دولت کی
منصفانہ گردش۔
\end{hal}
